% GENERATED FILE. Edit canonical lessons, then run npm run generate:books.
% canonical-source-hash: fnv1a64:b0a04a7a34f8bd46
% canonical-lessons: JA-C209-haru, JA-C209-aki, JA-C209-kisetsu, JA-C209-sakura, JA-C209-chiru, JA-R209-review-204-209

\chapter{Spring, Autumn and Cherry Blossom}
\label{ch:ja-spring-autumn-and-cherry-blossom}
\hlchaptermodality{209}

\begin{chapteropening}
\textbf{By the end of this chapter:} \emph{I can talk about the seasons ahead.}

Complete the last lesson of chapter 209: I can talk about the seasons ahead.
\end{chapteropening}

\section[\texorpdfstring{\ja{はる} (haru)}{はる (haru)}]{\ja{はる} (haru) — spring (the season)}

\label{lesson:JA-C209-haru}



\begin{quote}
Before the new one: say the Japanese for the sky, then the Japanese for a typhoon.
\end{quote}

\subsection*{You'll want to know: \ja{はる}}
\textbf{\ja{はる}} — \emph{haru} — \textquotedblleft{}spring (the season)\textquotedblright{}.

\begin{quote}
\textbf{\ja{らいげつは} \ja{はるです。}} — \emph{raigetsu wa haru desu} — \textquotedblleft{}Next month is spring.\textquotedblright{}
\end{quote}

\begin{cousinweb}
\textbf{\ja{はる}} (\emph{haru}) is a native Japanese word (wago), inherited rather than borrowed from Chinese; its kanji was attached later as a kun'yomi, and its older origin is not settled (JMdict, KANJIDIC2).
\end{cousinweb}

\begin{grammarlens}[title={What you've built}]
One more word for everyday life at A2.
\end{grammarlens}

\subsection*{Your turn}
\begin{itemize}
\raggedright
  \item \emph{Say it:} \emph{haru}
  \item \emph{Say it:} \emph{haru}, once more
  \item \emph{Say it:} \emph{raigetsu wa haru desu} — Next month is spring.
  \item \emph{Recall:} say the Japanese for the sky, then the Japanese for a typhoon, then say \emph{haru} again
\end{itemize}

\begin{tcolorbox}[breakable,colback=teal!4,colframe=teal!35!black,title={Before you move on}]
What does \ja{はる} mean? (\textquotedblleft{}Spring (the season)\textquotedblright{}.) Say it once more.
\end{tcolorbox}
\section[\texorpdfstring{\ja{あき} (aki)}{あき (aki)}]{\ja{あき} (aki) — autumn}

\label{lesson:JA-C209-aki}



\begin{quote}
Before the new one: say the Japanese for a typhoon, then the Japanese for spring.
\end{quote}

\subsection*{You'll want to know: \ja{あき}}
\textbf{\ja{あき}} — \emph{aki} — \textquotedblleft{}autumn\textquotedblright{}.

\begin{quote}
\textbf{\ja{らいねんの} \ja{あきに} \ja{りょこうに} \ja{いきます。}} — \emph{rainen no aki ni ryokō ni ikimasu} — \textquotedblleft{}I will go on a trip next autumn.\textquotedblright{}
\end{quote}

\begin{cousinweb}
\textbf{\ja{あき}} (\emph{aki}) is a native Japanese word (wago), inherited rather than borrowed from Chinese; its kanji was attached later as a kun'yomi, and its older origin is not settled (JMdict, KANJIDIC2).
\end{cousinweb}

\begin{grammarlens}[title={What you've built}]
One more word for everyday life at A2.
\end{grammarlens}

\subsection*{Your turn}
\begin{itemize}
\raggedright
  \item \emph{Say it:} \emph{aki}
  \item \emph{Say it:} \emph{aki}, once more
  \item \emph{Say it:} \emph{rainen no aki ni ryokō ni ikimasu} — I will go on a trip next autumn.
  \item \emph{Recall:} say the Japanese for a typhoon, then the Japanese for spring, then say \emph{aki} again
\end{itemize}

\begin{tcolorbox}[breakable,colback=teal!4,colframe=teal!35!black,title={Before you move on}]
What does \ja{あき} mean? (\textquotedblleft{}Autumn\textquotedblright{}.) Say it once more.
\end{tcolorbox}
\section[\texorpdfstring{\ja{きせつ} (kisetsu)}{きせつ (kisetsu)}]{\ja{きせつ} (kisetsu) — a season}

\label{lesson:JA-C209-kisetsu}



\begin{quote}
Before the new one: say the Japanese for spring, then the Japanese for autumn.
\end{quote}

\subsection*{You'll want to know: \ja{きせつ}}
\textbf{\ja{きせつ}} — \emph{kisetsu} — \textquotedblleft{}a season\textquotedblright{}.

\begin{quote}
\textbf{\ja{はるは} \ja{いい} \ja{きせつです。}} — \emph{haru wa ii kisetsu desu} — \textquotedblleft{}Spring is a good season.\textquotedblright{}
\end{quote}

\begin{cousinweb}
\textbf{\ja{きせつ}} (\emph{kisetsu}) is Sino-Japanese (kango): two Chinese characters read with their on'yomi, readings borrowed from Chinese. They are \emph{ki}, \textquotedblleft{}season\textquotedblright{} (Middle Chinese \emph{kjwijH}) and \emph{setsu}, \textquotedblleft{}joint, period\textquotedblright{} (\emph{tset}) (KANJIDIC2; Middle Chinese in Baxter's transcription).
\end{cousinweb}

\begin{grammarlens}[title={What you've built}]
One more word for everyday life at A2.
\end{grammarlens}

\subsection*{Your turn}
\begin{itemize}
\raggedright
  \item \emph{Say it:} \emph{kisetsu}
  \item \emph{Say it:} \emph{kisetsu}, once more
  \item \emph{Say it:} \emph{haru wa ii kisetsu desu} — Spring is a good season.
  \item \emph{Recall:} say the Japanese for spring, then the Japanese for autumn, then say \emph{kisetsu} again
\end{itemize}

\begin{tcolorbox}[breakable,colback=teal!4,colframe=teal!35!black,title={Before you move on}]
What does \ja{きせつ} mean? (\textquotedblleft{}A season\textquotedblright{}.) Say it once more.
\end{tcolorbox}
\section[\texorpdfstring{\ja{さくら} (sakura)}{さくら (sakura)}]{\ja{さくら} (sakura) — a cherry tree; cherry blossom}

\label{lesson:JA-C209-sakura}



\begin{quote}
Before the new one: say the Japanese for autumn, then the Japanese for a season.
\end{quote}

\subsection*{You'll want to know: \ja{さくら}}
\textbf{\ja{さくら}} — \emph{sakura} — \textquotedblleft{}a cherry tree; cherry blossom\textquotedblright{}.

\begin{quote}
\textbf{\ja{らいしゅう} \ja{さくらが} \ja{さきます。}} — \emph{raishū sakura ga sakimasu} — \textquotedblleft{}The cherry trees will bloom next week.\textquotedblright{}
\end{quote}

\begin{cousinweb}
\textbf{\ja{さくら}} (\emph{sakura}) is a native Japanese word (wago), inherited rather than borrowed from Chinese; its kanji was attached later as a kun'yomi, and its older origin is not settled (JMdict, KANJIDIC2).
\end{cousinweb}

\begin{grammarlens}[title={What you've built}]
One more word for everyday life at A2.
\end{grammarlens}

\subsection*{Your turn}
\begin{itemize}
\raggedright
  \item \emph{Say it:} \emph{sakura}
  \item \emph{Say it:} \emph{sakura}, once more
  \item \emph{Say it:} \emph{raishū sakura ga sakimasu} — The cherry trees will bloom next week.
  \item \emph{Recall:} say the Japanese for autumn, then the Japanese for a season, then say \emph{sakura} again
\end{itemize}

\begin{tcolorbox}[breakable,colback=teal!4,colframe=teal!35!black,title={Before you move on}]
What does \ja{さくら} mean? (\textquotedblleft{}A cherry tree; cherry blossom\textquotedblright{}.) Say it once more.
\end{tcolorbox}
\section[\texorpdfstring{\ja{ちる} (chiru)}{ちる (chiru)}]{\ja{ちる} (chiru) — to scatter; to fall (of petals or leaves)}

\label{lesson:JA-C209-chiru}



\begin{quote}
Before the new one: say the Japanese for a season, then the Japanese for a cherry tree.
\end{quote}

\subsection*{You'll want to know: \ja{ちる}}
\textbf{\ja{ちる}} — \emph{chiru} — \textquotedblleft{}to scatter; to fall (of petals or leaves)\textquotedblright{}.

\begin{quote}
\textbf{\ja{あした} \ja{さくらが} \ja{ちります。}} — \emph{ashita sakura ga chirimasu} — \textquotedblleft{}The cherry blossoms will fall tomorrow.\textquotedblright{}
\end{quote}

Polite \textbf{\ja{ちります}}, \emph{chirimasu}; a godan verb.

\begin{cousinweb}
\textbf{\ja{ちる}} (\emph{chiru}) is a native Japanese word (wago), inherited rather than borrowed from Chinese; its kanji was attached later as a kun'yomi, and its older origin is not settled (JMdict, KANJIDIC2).
\end{cousinweb}

\begin{grammarlens}[title={What you've built}]
One more word for everyday life at A2.
\end{grammarlens}

\subsection*{Your turn}
\begin{itemize}
\raggedright
  \item \emph{Say it:} \emph{chiru}
  \item \emph{Say it:} \emph{chiru}, once more
  \item \emph{Say it:} \emph{ashita sakura ga chirimasu} — The cherry blossoms will fall tomorrow.
  \item \emph{Recall:} say the Japanese for a season, then the Japanese for a cherry tree, then say \emph{chiru} again
\end{itemize}

\begin{tcolorbox}[breakable,colback=teal!4,colframe=teal!35!black,title={Before you move on}]
What does \ja{ちる} mean? (\textquotedblleft{}To scatter; to fall (of petals or leaves)\textquotedblright{}.) Say it once more.
\end{tcolorbox}
\section[(dialogue)]{Words from chapters 204–209 — a review}

\label{lesson:JA-R209-review-204-209}



\begin{quote}
Say the Japanese for a trip and a Japanese inn.
\end{quote}

\subsection*{Your turn}
Say these aloud:

\begin{itemize}
\raggedright
  \item a trip, a Japanese inn, sightseeing, to drop in, lodging
  \item a ticket window, a station attendant, a ticket gate, one way, a round trip
  \item a vehicle, traffic, a traffic jam, a passenger, a booking
  \item the weather, to fall, snow, to clear up, to cloud over
  \item fine weather, cloudy weather, a weather forecast, the sky, a typhoon
  \item spring, autumn, a season, a cherry tree, to scatter
\end{itemize}

\begin{tcolorbox}[breakable,colback=teal!4,colframe=teal!35!black,title={Before you move on}]
A trip? (\textbf{\ja{りょこう}}, \emph{ryokō}.) A Japanese inn? (\textbf{\ja{りょかん}}, \emph{ryokan}.) Sightseeing? (\textbf{\ja{けんぶつ}}, \emph{kenbutsu}.) To drop in? (\textbf{\ja{よる}}, \emph{yoru}.) Lodging? (\textbf{\ja{やど}}, \emph{yado}.) A ticket window? (\textbf{\ja{まどぐち}}, \emph{madoguchi}.) A station attendant? (\textbf{\ja{えきいん}}, \emph{ekiin}.) A ticket gate? (\textbf{\ja{かいさつぐち}}, \emph{kaisatsuguchi}.) One way? (\textbf{\ja{かたみち}}, \emph{katamichi}.) A round trip? (\textbf{\ja{おうふく}}, \emph{ōfuku}.) A vehicle? (\textbf{\ja{のりもの}}, \emph{norimono}.) Traffic? (\textbf{\ja{こうつう}}, \emph{kōtsū}.) A traffic jam? (\textbf{\ja{じゅうたい}}, \emph{jūtai}.) A passenger? (\textbf{\ja{じょうきゃく}}, \emph{jōkyaku}.) A booking? (\textbf{\ja{よやく}}, \emph{yoyaku}.) The weather? (\textbf{\ja{てんき}}, \emph{tenki}.) To fall? (\textbf{\ja{ふる}}, \emph{furu}.) Snow? (\textbf{\ja{ゆき}}, \emph{yuki}.) To clear up? (\textbf{\ja{はれる}}, \emph{hareru}.) To cloud over? (\textbf{\ja{くもる}}, \emph{kumoru}.) Fine weather? (\textbf{\ja{はれ}}, \emph{hare}.) Cloudy weather? (\textbf{\ja{くもり}}, \emph{kumori}.) A weather forecast? (\textbf{\ja{てんきよほう}}, \emph{tenkiyohō}.) The sky? (\textbf{\ja{そら}}, \emph{sora}.) A typhoon? (\textbf{\ja{たいふう}}, \emph{taifū}.) Spring? (\textbf{\ja{はる}}, \emph{haru}.) Autumn? (\textbf{\ja{あき}}, \emph{aki}.) A season? (\textbf{\ja{きせつ}}, \emph{kisetsu}.) A cherry tree? (\textbf{\ja{さくら}}, \emph{sakura}.) To scatter? (\textbf{\ja{ちる}}, \emph{chiru}.)
\end{tcolorbox}
